\subsection{결과 논의 및 대응 방향}
\label{sec:discussion}

\subsubsection{SW 취약성 분석과 T1 물리 검증의 역할}

본 연구의 SW 에뮬레이션과 물리 실험은 서로 다른 연구 질문에 답한다.
SW 에뮬레이션은 목표 오류 효과가 발생했을 때 어떤 대수적 관계와
보안 영향이 형성되는지를 검증한다. 따라서 T2, RB 및 상류 오류
지점에 대한 E2 결과는 해당 지점의 취약성을 식별하고 분류하는 본
논문의 목적에 부합한다. 다만 이 결과를 물리 글리치의 발생률로
해석하지 않는다.

T1 실험은 이러한 의미론적 관계가 실제 클럭 글리치로 생성될 수 있는지
한 단계 더 검증한다. 세 곱셈 반복에 대응하는 글리치 구간을 확인하고,
집중 실행과 별도 실행에서 각각 $\svec_1$ 768/768을 복원하였다.
이는 유용한 결함 응답 두 개만을 선택한 사후 결과가 아니라, 탐색과
재현 과정을 함께 기록한 물리 실현성 증거이다. 동시에 최소 악용 출력
집합과 전체 측정 비용을 분리하여 공격의 조건을 명확히 한다.

\subsubsection{구현 구조가 오류 효과에 미치는 영향}

T1 인과대조 구현에서는 곱 결과 버퍼의 초기 상태가 물리 오류 결과에
직접 영향을 주었다. 버퍼가 0으로 초기화된 경우 스킵된 반복은 명확한
항 제거로 나타났으며, 다항식 블록별 비밀 성분을 복원할 수 있었다.
반면 이전 값이 남아 있는 구현에서는 동일한 명령어 생략이 stale value를
응답에 반영할 수 있다. 따라서 소스 코드의 수식뿐 아니라 컴파일된
반복 구조와 버퍼 생명주기를 함께 분석해야 한다.

T2에서는 SW 훅으로 벡터 전체의 nonce 항을 제거할 수 있었지만,
미수정 융합 구현의 단일 글리치는 동일한 효과를 만들지 못하였다.
이는 구현 의존성을 실험적으로 시사하지만, 특정 융합 구조가 영구적인
방어를 제공한다는 의미는 아니다. 컴파일러 최적화, 루프 분할, 다른
마이크로아키텍처 또는 다중 글리치에서는 오류의 범위가 달라질 수 있다.

\subsubsection{공격 범위와 후속 고도화}

물리 T1의 복원과 판정에는 연구용 내부 응답과 참 $\svec_1$ 스트림을
사용하였다. 이는 클럭 글리치와 블록 제거 사이의 인과관계를 확인하기
위한 계측이며, 표준 직렬화 서명만을 사용하는 종단 공격을 주장하기
위한 자료는 아니다. 또한 복원 대상은 전체 서명키 패키지가 아니라
비밀 성분 $\svec_1$이다.

따라서 signature-only 복원, 실제 복원 후보의 공개키 일관성 검사 및
표준 verifier를 통과하는 위조는 현재 취약성 분석의 필수 전제가 아니라
후속 공격 고도화 과제이다. 이 과정을 추가하면 E3 물리 인과 검증에서
E4--E5 공격 영향 검증으로 범위를 확장할 수 있다.

RB 역시 T1/T2의 직접 항 제거와 구분해야 한다. SW 실험은 거부되어야
할 응답이 후속 경로로 진행될 수 있음을 보였지만, 키 복구나 정보량을
측정하지 않았다. 따라서 RB는 다수 출력에 대한 likelihood 또는
key-ranking 분석이 완료되기 전까지 잠재적 분포 취약성으로 한정한다.

\subsubsection{대응 방향}

본 논문의 주기여는 취약성 분석과 T1 물리 검증이며, 대응기법은 분석에서
도출되는 최소 보안 요구사항으로만 제시한다.

\begin{table*}[t]
\centering
\caption{Security requirements derived from the vulnerability analysis}
\label{tab:defense-requirement}
\small
\begin{tabularx}{0.94\textwidth}{@{}l X X@{}}
\toprule
Threat & Required property & Possible direction \\
\midrule
T1/T2 항 제거 &
응답 계산과 중간 버퍼의 데이터 흐름 무결성 &
독립 재계산, 중간값 태그 및 바이너리 수준 자체검사 \\
RB &
서명 경계와 거부 상태의 무결성 &
$B_0/B_1$ 재검사 및 거부 loop 상태 보호 \\
검사 분기 생략 &
단일 출력 분기에 의존하지 않는 오류 반응 &
감염형 출력 변환 또는 중복 제어흐름 \\
\bottomrule
\end{tabularx}
\end{table*}

응답 잔차를 최종 출력에 확산하는 감염형 처리는 가능한 방향 중 하나이다
~\cite{yen2000,seo2013}. 저장소의 IRV 프로토타입도 이러한 방향을
탐색하지만, 전체 검사 경로의 무분기성, 다중 오류 내성 및 공정한 비용
비교는 추가 검증이 필요하다. 따라서 본 논문에서는 IRV를 완성된 핵심
대응기법으로 주장하지 않고 후속 연구로 남긴다.

\subsubsection{연구의 한계}

물리 실험은 STM32F405, 특정 컴파일러와 펌웨어 설정 및 클럭 글리치에
한정된다. T1 결과는 버퍼 사전 초기화와 고정 nonce를 사용하는
인과대조 구현에서 얻었고, T2의 0/650 결과는 물리적 불가능성을
증명하지 않는다. SW 결과는 의미론적 취약성을 검증하지만 각 지점의
물리 발생률을 제공하지 않는다. 이러한 범위는 주장과 표에서 분리하여
표시하였으며, 다른 구현과 주입 수단에 대한 평가는 후속 연구로 둔다.
